\scanpage{37}

\uline{L'idée de} W-M. \deleted{serait} \inserted{\uline{semble pouvoir s'exprimer}} \uline{plus ou moins} \inserted{\uline{en disant}} \uline{que ce module est un invariant qui, à isom canonique près, ne dépend pas du choix du relèvement $X$ de $X_0$, et qui serait donc un invariant intrinsèquement attaché à $X_0$, qui serait alors par définition la cohomologie} \deleted{de D.R.} \uline{``$p$-adique'' \unsure{de celui-ci}.}

\margin{``\unsure{Some questions}'' \unsure{Journal Indien} 63/64\ldots}

\uline{Remarques} a) Cette construction marcherait aussi, mutatis mutandis, si on ne peut relever $X_0$ qu'en un \uline{schéma formel} lisse sur $S$.

b) D'après le texte publié de W-M, \deleted{il me semble qu'il} on pourrait \deleted{croire} être amené à penser que seul
\[
H^i_{\mathrm{DR}}(X/S)\otimes_W K
\]
\inserted{$K$ corps des fractions de $W$}
\deleted{\unsure{on}} devrait avoir la propriété d'invariance en question. De plus, la construction et la propriété invariance n'y sont traités que pour $X_0$ \uline{affine}

\margin{36}
